% SOURCE_STATEMENT: "For Reynolds-averaged equations modeling the variances of liquid-water potential temperature and total-water mixing ratio in a precipitating cloudy atmospheric boundary layer, can the equations represent explicitly the slow positive feedback that produces the variance without an additional stochastic submodel?"
% Preparation and final self-audit: CONTRIBUTING.md prompts C and D; source/audit artifacts retained in the local physics_statements_2026_10_11 batch.
% !TeX program = lualatex
\documentclass[11pt]{article}
\usepackage[margin=25mm]{geometry}
\usepackage{fontspec}
\setmainfont{Latin Modern Roman}
\usepackage{amsmath,amssymb,mathtools}
\usepackage{unicode-math}
\setmathfont{Latin Modern Math}
\usepackage{xurl,hyperref}
\hypersetup{colorlinks=true,urlcolor=blue}
\setcounter{secnumdepth}{0}
\setlength{\parindent}{0pt}
\setlength{\parskip}{5pt}
\setlength{\emergencystretch}{3em}
\newcommand{\sref}[2]{\hyperlink{source:#1:#2}{[S#2]}}
\newcommand{\cataloguescope}[1]{\paragraph{Recorded target.}#1}
\newcommand{\R}{\mathbb R}
\newcommand{\C}{\mathbb C}
\newcommand{\N}{\mathbb N}
\newcommand{\Z}{\mathbb Z}
\newcommand{\Q}{\mathbb Q}
\newcommand{\F}{\mathbb F}
\newcommand{\eps}{\varepsilon}
\begin{document}
\section{Variance Feedback in Cloud Boundary-Layer Models}\label{statement}

\subsection{Definitions and mathematical statement}

\paragraph{Shared notation.}$\mathbb N=\{1,2,\ldots\}$; the standard number fields have their usual meanings. All additional conventions are stated locally.


\paragraph{Definitions and conventions.}
A Reynolds decomposition writes a fluctuating quantity as its mean plus a fluctuation, and variance is the mean square of that fluctuation. Liquid-water potential temperature and total-water mixing ratio are the two cloud variables in the source closure equations. Their variance budgets contain production, transport and dissipation terms. A slow positive feedback is a coupling in which a change in a mean field enhances variance production, which in turn reinforces that change.

\paragraph{Target retained from the canonical source catalogue.}
\paragraph{Catalogue wording retained for traceability.}
For Reynolds-averaged equations modeling the variances of liquid-water potential temperature and total-water mixing ratio in a precipitating cloudy atmospheric boundary layer, can the equations represent explicitly the slow positive feedback that produces the variance without an additional stochastic submodel?

\paragraph{Source scope and qualification.}
The report\textquotesingle s precipitating boundary-layer budget and feedback mechanism are the target; the extracted statement does not supply the full closure equations. Adding an independent stochastic production model would not meet the stated without-an-additional-submodel requirement.
\subsection{Short English statement}
Can deterministic Reynolds-averaged cloud equations explicitly generate the observed slow feedback that raises thermodynamic variances?
\subsection{Sources}
\begin{itemize}
\item[\hypertarget{source:statement:1}{S1}] ULAM.AI and the UnsolvedMath Contributors, UnsolvedMath, catalogue entry. Catalogue curation under CC BY 4.0; definitions and English exposition added locally.
\url{https://huggingface.co/datasets/ulamai/UnsolvedMath}.
\item[\hypertarget{source:statement:2}{S2}] Moist Processes in the Atmosphere (2019), Oberwolfach Reports, DOI 10.4171/owr/2019/7, pp. 465--512. Primary report accessed 2026-10-11; relevant PDF pages 14.
\url{https://doi.org/10.4171/owr/2019/7}.
\item[\hypertarget{source:statement:3}{S3}] Stochastic parameterization review (catalogue-listed context; not independently consulted in this preparation).
\url{https://journals.ametsoc.org/view/journals/bams/98/3/bams-d-15-00268.1.xml}.
\end{itemize}
\label{end:statement}
\end{document}
